\section{Joint interpretation}

The empirical contribution is a common prescription that connects a
galactic acceleration scale, a retained mass distribution, tidal evolution,
and cosmological energy transfer. The calculations use the same
microscopic entropy and conversion fraction in several settings. Their
success would be more constraining than unrelated curve fits because
changing those coefficients would move several observables at once.

The strongest present numerical agreements are the LMC enclosed masses,
the two coefficients of the published acceleration relation under the
stated linear projection, and the near-equal cosmological likelihood
with a tied abundance. Crater II supplies a concrete dynamical mechanism
and a favorable comparison with a specified NFW population. Its evidence
depends strongly on the adopted error budget. The tidal-dwarf calculation
tests inheritance at the level of a small dynamical-mass sample. Antlia 2
and the independently measured cosmological mixing normalization remain
active constraints on the model.

\subsection{The derivation has identifiable physical commitments}

\begin{table}[htbp]
\small
\begin{tabularx}{\textwidth}{P{.24\textwidth}YY}
\toprule
Step & Result supplied & Physical commitment or outstanding test \\
\midrule
Finite face and cell ensemble
& $16=7+9$; 1680 cell states; $g=7.4198000$ and $\eta_*=0.02986684$.
& The spin, matching vertex, occupancy rule, closure invariant, and
evidence measure define the ensemble. \\
Horizon and transverse response
& $\epsilon=g/(2\pi)^2$; $a_0=\epsilon cH$; Bose response.
& Phase loading, horizon contact, and the mode-energy mapping connect
the count to acceleration. \\
Retained galactic stock
& An inertial positive mass component, inherited by carriers and
transported by collisionless evolution.
& Assembly budgets, carrier assignment, profile preparation, and
relaxation must be specified independently of the target observables. \\
Cosmological conversion
& $\eta=\eta_*$ and an initial abundance near $\epsilon^{-1}$.
& Common energy conversion, the selected pre-acoustic formation map
(whose microscopic vertex remains open), and the absolute vacuum
matching condition enter the physical identification. \\
Assembly and mixing
& A cumulative release function and a weighted vacuum integral.
& Retention gates, progenitor labels, and transport set the normalization;
the current expansion fit determines one mixing parameter. \\
Geometry and perturbations
& Reduced Newtonian equations and no-slip identity proved from the
specified action; carrier equations proved on their regular branch.
& The continuum embedding, constrained director dynamics, and the
physical matching of the effective actions remain separate requirements. \\
\bottomrule
\end{tabularx}
\caption{The logical connection between the microscopic model and the
observables. An exact finite calculation can coexist with conditional
physical identifications and an incomplete continuum derivation.}
\end{table}

This chain explains why the empirical results bear on the microstructure:
the model does not choose a separate acceleration coefficient for the LMC
and the redshift sample, or a separate release fraction for each
cosmological data combination. It also identifies which steps the data
cannot determine by themselves. A fit to an effective gravitational field
does not uniquely identify the microscopic state space that produced it.

\subsection{Which conventional assumptions are challenged}

For the LMC, the comparison challenges the applicability of an
instantaneous external-field suppression to a previously established
gravitating stock. A retained distribution continues to contribute after
the ambient field changes. The tested QUMOND prescriptions and the stock
model make different mass predictions with the same stated baryons.
Joint stream and reflex modeling can test that distinction more directly
than comparing scalar enclosed-mass summaries.

For Crater II, the issue is the central structure and stripping response
of the excess gravitating mass. The tested capacity profile loses enough
support to produce a cold, extended remnant under the adopted orbit and
initial conditions. Standard-concentration NFW halos under the tested
abundance prior perform poorly in that comparison. Low-concentration
and cored dark-matter alternatives require their own matched tests.
Assembly history already matters in conventional dark-matter dynamics;
history dependence alone cannot distinguish the two frameworks.

For the acceleration scale, the proposed change is a cosmological
normalization $a_0\propto H$. Its observed value at one epoch would cease
to be an independent constant. The published intermediate-redshift
relation is consistent with the fixed reference curve under the stated
projection. A direct galaxy-level likelihood would test both its curvature
and its scatter.

For dark energy, the proposed change is a vacuum density generated by
irreversible source-identity loss. The microscopic conversion and
formation abundance restrict a model that still fits the CMB and distance
data. The remaining release normalization must be obtained independently
before the observed vacuum abundance can be called a prediction of the
microstructure alone. The current evidence supports these mechanisms as
testable alternatives; it does not establish the elimination of all
dark-matter models, all MOND formulations, or the absolute cosmological
constant problem.

\subsection{A joint test of the same construction}

The galactic comparisons use the reference
$H_0=67.4\,\mathrm{km\,s^{-1}\,Mpc^{-1}}$, while the matched linear-lensing
release fit gives $68.9389$. If horizon contact is evaluated at the latter
value, the local scale rises by about $2.3\%$. A joint fit must propagate
the fitted $H(z)$ into galaxy assembly and the acceleration relation,
including any finite relaxation time of an existing stock. The
reference-curve agreement and the cosmological best fit are presently
separate calculations.

The same requirement applies to the Milky Way potential. The available
moving-LMC calculations use a Galactic mass within 50 kpc about $17\%$
above the Brooks et al. inference in the stock model, and about $36\%$
above it in the NFW model. A joint Milky Way--LMC fit would alter the
satellite pericentres and the stripping histories. Those changes should
be propagated to both Crater II and Antlia 2 before assigning a combined
evidence ratio. \source{G4}\cite{brooks}

A decisive program follows from these dependencies:
\begin{enumerate}
\item Infer the Milky Way and LMC jointly using streams, halo reflex
kinematics, and a common source-conserving force model. Predict additional
stream observables from the resulting mass distributions.
\item Evolve the dwarf populations with resolved self-gravity and
prespecified progenitor priors. Apply the actual stellar selection,
velocity-gradient, and membership inference to their mock observations.
Include concentration scatter and physically motivated cores in the
dark-matter comparison.
\item Fit the intermediate-redshift galaxy data directly with the
cosmological $H(z)$ from the same model. Preserve correlations between
baryonic masses, pressure support, geometry, and acceleration.
\item Measure retention and source-label loss in converged hydrodynamic
runs with a label prescription fixed in advance. Use that release history
with the existing positive-allocation perturbation implementation and
a matched full cosmological likelihood.
\end{enumerate}
Each step tests an existing physical commitment of the construction.

\appendix
\section{Provenance and numerical verification}

\subsection{Scope of the verification}

The Lean 4.35.0 library contains 154 modules and 2,812 audited theorem
constants, including its legacy sublibrary. The manuscript ledger records
874 statements as proved, 90 in part, 153 as physical premises, 64 as
data comparisons, and 19 as simulation outputs. Its other entries
distinguish numerical calculations, definitions, and statements awaiting
formalization. The saved audit reports only the standard Lean axioms
and no admitted proofs. The counts include mathematical helper results.
\source{T4}

For this paper, the relevant proofs cover the finite ensemble and closure
root, the restricted freedom in the marked-register identification,
positive finite allocation, and radial relaxation with preserved angular
marginals. First variation with compactly supported tests gives the
reduced two-potential and Newton equations. Decay at infinity gives
$\Phi=\Psi$. The carrier-action proof assumes a smooth interpolating
function; the adopted singular derivative at zero baryonic field limits
its direct classical application. The constrained director equation,
full cosmological perturbations, and observational likelihoods retain
their separate analytic or numerical status.

The numerical review independently evaluates the finite ensemble,
redshift projection, LMC mass profile and QUMOND flux, archived Crater II
likelihoods, and a reference release history. The current verification
record uses the unrounded acceleration coefficient and imports the
higher-accuracy matched cosmology endpoints. The associated figures
use those records. The archived release-history and hydrodynamic panels
retain their original linear-depletion calibration.

The dated research records include the exponential-allocation cosmological
optimization, finite-work and fluid tests, merger controls, observed cluster
reductions, and outlier audits. They were not rerun to typeset this paper.
Six extension modules contribute 119 explicit theorem declarations,
independently checked with the available Lean/mathlib 4.19 runtime and
added to the pinned library's import graph. Their build and axiom records
are separate from the upstream 4.35 audit. \source{T5} Source inspection and saved
audits establish which calculation supports each statement. Numerical
convergence, selection effects, and physical matching are assessed with
the relevant empirical result. SHA-256 records preserve the identities
of the source and result snapshots.

\clearpage
\subsection{Local source key}

Paths below are relative to the supplied \path{Sol and Fable} workspace.
The source package includes the document, figures, numerical review, and
verification script; the original simulation archives remain in that
workspace.

\small
\begin{longtable}{P{.065\textwidth}P{.88\textwidth}}
\toprule
Key & Source and use \\
\midrule
T1 & \path{01 Current Papers/Astra Full/Entropic Scalar EFT.tex}. Current microscopic
construction, effective galaxy dynamics, release accounting, and
covariant closure; especially sections 5--6, 15, 20, N.8, and O.3--O.15.
The numerical review records its source snapshot by hash. \\
T2 & \path{02 Research/Spacetime consolidation 2026-10-16/MICROSTRUCTURE_TO_SPACETIME.md}.
Finite geometry, interaction, matter, and clock
construction. \\
T3 & \path{01 Current Papers/Microstructure CDT/MICROSTRUCTURE_CDT.tex}.
Selected CDT realization, normalization conventions, and finite
geometry--matter--clock tests, with conventions aligned to the current
Entropic Scalar EFT manuscript. \\
T4 & \href{https://github.com/jakechinitz/Entropic-EFT-Lean/tree/432da8c4a8b169a17146b53e4155d6b37915a12f}{Entropic-EFT-Lean, 23 October 2026 snapshot}.
Local library: \path{02 Research/Unified closure 2026-10-21/formalization}.
Claim ledger, theorem sources, and saved axiom/citation audits. \\
G1 & \path{04 Notes and References/Technical Notes/history_gravity_reproduction_2026-09-27.md}.
Analytic galaxy checks and reproducible provenance for the stock model. \\
G2 & \path{02 Research/LMC and satellites checks/qumond_flux.py} and
\path{02 Research/LMC and satellites checks/output.txt}.
External-field QUMOND surface-flux comparison. \\
G3 & \path{02 Research/Crater II and Antlia 2 vs LCDM/PREREGISTRATION.md},
\path{PREREGISTRATION.sha256}, and \path{analyse_grid.py} in that folder.
Core observation likelihood, priors, and hash identity. \\
G4 & \path{02 Research/Crater II and Antlia 2 vs LCDM/RESULTS_with_LMC.md},
\path{analyse_lmc.py}, and \path{lmc_results.zip} in that folder.
Current moving-LMC trajectories and evidence. \\
G5 & \path{04 Notes and References/Technical Notes/carrier_hierarchy_calcs_2026-09-26.md}.
The active-carrier inheritance calculation for tidal dwarfs. \\
G6 & \path{02 Research/Crater II and Antlia 2 vs LCDM/bf_robustness.txt} and
\path{antlia2_ji_mock/mock_approx_analysis.txt} in that folder.
Model-error sensitivity and approximate observational mocks. \\
G7 & \path{02 Research/Stock adjustment resolution 2026-10-21/START HERE.txt};
\path{02 Research/Halo replication audit 2026-10-21}.
Radial adjustment, retained angular marginals, and fixed-target halo tests. \\
C1 & \path{02 Research/Full Planck tied abundance/model.py},
\path{RESULTS.md}, and \path{fits_lenscorr.jsonl} in that folder;
\path{02 Research/Dark energy prediction tests/common.py}.
Matched likelihood, data handling, Boltzmann mapping, and reference
\LCDM\ result. \\
C2 & \path{02 Research/Release zoom/RESULTS.md}, \path{zoom_fits.jsonl},
\path{zoom_fit.py}, \path{o12_release.py}, \path{Frel_table.npy},
\path{FZ_ref.npy}, and \path{kmix_summary.txt} in that folder.
Archived PPF release fit, reference history, and hydrodynamic mixing test. \\
C3 & \path{02 Research/Formation operator check/RESULTS.md}.
Formation abundance and the $5.333$ mass convention. \\
C4 & \path{02 Research/High leverage closure 2026-10-21/fit_comparison.json},
\path{fit_physical_geodesic/result.json}, \path{fit_physical_lcdm/result.json},
\path{cosmology_checks.json}, and \path{growth_independent_check.json} in that folder.
Direct clock-fluid implementation, matched linear-lensing fits, and numerical checks. \\

T5 & \path{02 Research/Manuscript cluster and dynamics integration 2026-10-23}.
Preserved source hashes, integration diffs, six-module manifest, compilation
checks and PDF validation. The extension sources are also under
\path{Entropic/Extensions20261023} in the Lean library. \\
U1 & \path{02 Research/Unified empirical program 2026-10-23/execution_register.json}.
Fifteen-sector evidence register; S01--S15 identify the selected datasets,
saved results, source URLs, and remaining tests. \\
C5 & \path{02 Research/Unified empirical program 2026-10-23/results/cosmology_final.json}.
Exponential-allocation fit, matched reference, and saved observables. \\
D1 & \path{02 Research/Connected theory closure 2026-10-23}.
Depletion action, unchanged linear domain, lag and noise bounds, and
conditional tensor calculation; 42 Lean declarations. \\
D2 & \path{02 Research/Effective dynamics and coupling 2026-10-23}.
Finite-work contact, warm fluid stress, compression and feedback checks;
36 Lean declarations. \\
CL1 & \path{02 Research/Conservative cluster merger 2026-10-23}.
Seventeen encounter/control runs, conservation and refinement records;
8 Lean declarations. \\
CL2 & \path{02 Research/Observed cluster tests 2026-10-23}.
X-COP and MENeaCS profiles, Bullet v1 targets, source hashes, and
15 numerical/data-handling checks. \\
CL3 & \path{02 Research/Cluster outlier closure 2026-10-23}.
Conditional pressure and transport audit, A644 positivity test, twelve
finite transport maps, 14 saved checks, and 22 Lean declarations. \\
\bottomrule
\end{longtable}
\normalsize

\subsection{Reproduction commands}

With Python, NumPy, SciPy, and Matplotlib installed, run
\begin{quote}\ttfamily
python reproduce\_review.py -{}-workspace "PATH/TO/Sol and Fable"
\end{quote}
The script evaluates archived inputs without launching simulation jobs.
The supplementary \path{refresh_verified_results.py} updates the
unrounded redshift arithmetic and imports the saved exponential-allocation
fit. It preserves prior outputs before writing. Its
\texttt{-{}-legacy-linear} option reproduces the earlier matched benchmark.
The broader calculations and source records are identified by U1 and C5.
Build the publication from
the workspace root with
\begin{quote}\ttfamily
python .organization/build\_pdfs.py -{}-paper empirical
\end{quote}
The build preserves previous PDFs and stores compiler files in the
workspace build directory. The source archive supplies the figures for
standalone typesetting.

The public dark-energy release and DES data repository are linked in the
references. The current physical mixing extension is identified by its
local archive and hashes; availability of the public reproduction bundle
does not by itself make every subsequent calculation public.

\section{Notation}
\small
\begin{tabularx}{\textwidth}{P{.23\textwidth}Y}
\toprule Symbol & Meaning \\
\midrule
$g$ & Cell sharing entropy in natural-log units. \\
$\epsilon=g/(4\pi^2)$ & Transverse response loading. \\
$\eta_*$ & Closure evidence coefficient; also the release fraction under
the common energy-conversion prescription. \\
$a_0$ & Characteristic galactic acceleration scale. \\
$g_b$ or $g_N$ & Newtonian acceleration produced by the assigned baryonic
source. \\
$n_B$ & Occupation of the adopted transverse mode. \\
$M_b$, $M_b^L$ & Present baryonic mass and ancestral Lagrangian baryonic
mass, respectively. \\
$C_R$ & Adopted retained-stock cap relative to present baryonic mass. \\
$R_h$ & Projected stellar half-light radius. \\
$\sigma_{\rm los}$ & Line-of-sight velocity dispersion, with sample and
gradient treatment specified for each comparison. \\
$Z$, $B$ & Prior-integrated likelihood and its ratio between two models. \\
$F_Z$ & Fraction that has first entered a baryon-retaining well. \\
$F_{\rm rel}$ & Cumulative fraction processed by the source-identity
release kernel. \\
$f$ & Mixing time in units of the formation halo dynamical time in the
pure-exponential fit. \\
$\rho_{c0}^{\rm wb}$ & Would-be present commitment density before release
losses, extrapolated as dust. \\
$I_{\rm rel}$ & Linear-benchmark release integral weighted by $a^{-3}$. \\
$I_{\eta,\rm rel}$ & Positive-allocation integral with the additional
survival factor $e^{-\eta_*F_{\rm rel}}$. \\
$\omega_i=\Omega_i h^2$ & Physical density parameter in the cosmological
fit. \\
\bottomrule
\end{tabularx}
\normalsize
